\documentclass[../thesis.tex]{subfiles}

\begin{document}
	\counterwithout{section}{chapter}
	\setcounter{section}{0}

	Описание разработанного в рамках проекта сервиса пакетной обработки данных.

	\section{Формулировка}

	Исследуется распределение величины продольного поля на оси многокатушечного соленоида.

	Внутри железного ярма установлено 5 концентрических токоведущих катушек. Вся конструкция имеет
	аксиальную симметрию. Необходимо найти распределение продольной компоненты магнитного поля вдоль оси
	системы и интеграл этой компоненты.

	Схематично, исследуемая область изображена на рисунке \ref{fig:4}:

	\begin{figure}[H]
		\centering
		\includegraphics[width=1\textwidth]{assets/img/sol/geo.png}
		\caption{Схематичное изображение многокатушечного соленоида}
		\label{fig:4}
	\end{figure}

	Для всей исследуемой области известны значения магнитной проницаемости материалов. Для катушек
	заданы количество витков и ток.
\end{document}